\documentclass[11pt,noanswers]{exam} %noanswers
\usepackage[usenames]{color} %used for font color
\usepackage{amssymb} %maths
\usepackage{amsmath} %maths
\usepackage[utf8]{inputenc} %useful to type directly diacritic characters
\usepackage{siunitx}
\usepackage{mathtools}
\usepackage[version=4]{mhchem}
\usepackage{tabto}
\newcommand{\qtty}[2]{\left(\qty{#1}{#2}\right)}
\sisetup{display-per-mode = symbol }
\sisetup{inter-unit-product = \ensuremath { { } \cdot { } } }
\usepackage[letterpaper, total={7in, 10in}]{geometry}
\usepackage{graphicx} % Required for inserting images
\renewcommand{\epsilon}{\varepsilon}
\renewcommand{\vec}[1]{\mathbf{#1}}
\newcommand{\mat}[1]{\begin{bmatrix*}[r] #1 \end{bmatrix*}}

\begin{document}
\flushleft
Name:\ \makebox[4in]{\hrulefill} \hspace{1in} Date:\ \makebox[1in]{\hrulefill} \\

\bigbreak
\textbf{Quiz:} \tabto{2cm} \S 1.4 The Matrix Equation $A\mathbf{x} = \mathbf{b}$ \\
\tabto{2cm} \S 1.5 Solution Sets of Linear Systems

\flushleft


\begin{questions}
    \question Multiply the following:
    \begin{equation*}
        \mat{2 & 8 & 5 \\ -8 & 2 & 6} \mat{2 \\ 1 \\ 3}
    \end{equation*}
    \begin{solution}
        \begin{equation*}
            \mat{-3 \\ 4}
        \end{equation*}
    \end{solution}
    
    \question Let $A = \mat{-4 & 8 \\ 1 & 5 \\ 2 & 7}$.
    \begin{parts}
        \part State a vector $\vec{b}$ such that the product $A\vec{b}$ is \emph{undefined}.
        \part Suppose $A$ is multiplied by a vector $\vec{v}$ such that the product \emph{defined}. The result will be a vector with how many entries?
    \end{parts}
    \begin{solution}
        \begin{parts}
            \part $\vec{b}$ needs to have exactly two entries.
            \part $A\vec{v}$ will be a vector with exactly three entries.
        \end{parts}
    \end{solution}

    \question Write the following system of equations as a vector equation involving a linear combination of vectors.
    \begin{align*}
        5x_1 - 2x_2 - x_3 &= 2 \\
        4x_2 + 3x_3 &= -1
    \end{align*}
    \begin{solution}
        \begin{equation*}
            x_1 \mat{5 \\ 0} + x_2 \mat{-2 \\ 4} + x_3 \mat{-1 \\ 3} = \mat{2 \\ -1}
        \end{equation*}
    \end{solution}
    
    \question Write the following system of equations as a product of a matrix and a vector.
    \begin{align*}
        3x_1 + x_2  - 2x_3 &= 4 \\
        5x_1 - 4x_3 &= -1
    \end{align*}
    \begin{solution}
        \[ \mat{3 & 1 & -2 \\ 5 & 0 & -4} \mat{4 \\ -1} \]
    \end{solution}

    \question Consider the ``system" of equations comprised of the single system $-2x_1 - 8 x_2 + 16x_3 = 0$ where $x_2$ and $x_3$ are considered to be free variables.
        \begin{parts}
            \part Describe the solution to this system as a linear combination of vectors 
            \part Describe the solution to this system geometrically.
        \end{parts}
    \begin{solution}
        \begin{parts}
        \part $\vec{x} = x_2 \mat{-4 \\ 1 \\ 0} + x_3 \mat{8 \\ 0 \\ 1}$
        \part It is a plane in $\mathbb{R}^3$.
        \end{parts}
    \end{solution}

    \question Find the general solution of the homogenous system, expressing your answer as a vector.
    \begin{align*}
        x_1 + 2x_2 - 3x_3 &= 0 \\
        4x_1 + 7x_2 - 9x_3 & = 0\\
        -x_1 -4x_2 + 9x_3 &= 0
    \end{align*}
    \begin{solution}
        $\vec{x} = \mat{x_1 \\ x_2 \\ x_3} = x_3 \mat{-3 \\ 3 \\ 1}$
    \end{solution}

        \question Consider the following matrix $A$ and vector $\vec{b}$:
        \begin{equation*}
            A = \mat{1 & 2 & -1 \\ 0 & 1 & 3 \\ -2 & -2 & k} \hspace{2cm} \vec{b} = \mat{-1 \\ 2 \\ b_3}
        \end{equation*}
        \begin{parts}
                \part For what value(s) of $k$ will $A\vec{x} = \vec{0}$ have a nontrivial solution. Find the non-trivial solution(s) in this (these) cases.
                \part For what value(s) of $b_3$ will $A\vec{x}=\vec{b}$ have a solution? Find it.
        \end{parts}
        \begin{solution}
        \begin{parts}
            \part $k=8$ yields the nontrivial solution $\vec{x}=\mat{7 \\ -3 \\ 1}$
            \part $b_3 = 6$ yields the solution $\vec{x} = \mat{7 \\ -3 \\ 1} + \mat{5 \\ -2 \\ 0}$
        \end{parts}
        \end{solution}
\end{questions}


\end{document}
